% !PRINTED 613
\weberpage{613}
%<<CONT
. If we take as our group the order $O' = [Q]$ with the conductor $Q$, which is defined by the congruence condition that every number in $O'$ shall be congruent to a rational number modulo $Q$, then we have a congruence group whose excludent and modulus is $= Q$. The classes correspond to the classes of forms of the discriminant $D = Q^2\Delta$. The second case includes the first within itself (for $Q = 1$). The class field is here the class field $\mathfrak{K}(D)$, which, to distinguish it from more general ones, we shall call the \gesperrt{order field}.

For if $(k)$ is a class invariant, then (\S\,122)
\[
(k)^p \equiv (k) \quad (\mathrm{mod}\ \mathfrak{P})
\]
is satisfied only when $p$ is representable by the principal class of the discriminant $D$. Then $p$ splits in $\Omega$ into two conjugate prime ideals $\mathfrak{p}$, $\mathfrak{p}'$, which are representable by \gesperrt{numbers} in $O$ or in $O'$, and hence are principal ideals.

The field $\mathfrak{K}(\Delta)$ is the \gesperrt{class field of $\Omega$ in the narrower sense}, the one that corresponds to the ideal classes of this field, and it may also be designated the \gesperrt{principal class field} $\mathfrak{K}(\Omega)$.

In every ideal class, both with respect to $O$ and with respect to $O'$, there are therefore infinitely many prime ideals of the first degree. In other words: if $\mathfrak{a}$ is an arbitrary ideal (coprime to $Q$), then there are infinitely many prime ideals $\mathfrak{p}$ which are equivalent to $\mathfrak{a}$ with respect to $O'$. Then, by \S\,95, the forms corresponding to the ideals $\mathfrak{a}$ and $\mathfrak{p}$ are also equivalent, and $p$ is \gesperrt{representable by the form} belonging to $\mathfrak{a}$. With this the theorem is proved:

1. \gesperrt{By every primitive form of the discriminant $D$ infinitely many primes are representable}\footnote{This theorem, for negative as well as for positive discriminants, was first proved by \gesperrt{Dirichlet}, but the proof was published only in one special case (Bericht der Berliner Akademie 1840, Werke Bd.\ I, S.\ 497). The detailed proof was given by H. \gesperrt{Weber} (Mathematische Annalen XX, 1882). A generalisation likewise originating with \gesperrt{Dirichlet}, to which we shall come back in the next section, was proved by A. \gesperrt{Meyer} (Crelles Journ., Bd.\ 103). The proof given here in the text rests on a different foundation, but for the present it of course applies (so long as the class field for positive discriminants is not known) only to negative discriminants.}.

\begin{center}
\S\,167. \textbf{Primes in Linear Forms.}
\end{center}

We now define a number group $A$ in $\Omega$ as follows:
%CONT>>
